% ------------------------------------------------------------------------------
% Christian Schramm
% ------------------------------------------------------------------------------
%
%
% Dokumentenkopf
% ------------------------------------------------------------------------------
\documentclass[
  11pt, % Schriftgröße
  DIV10,
  ngerman, % für Umlaute, Silbentrennung etc.
  a4paper, % Papierformat
  oneside, % einseitiges Dokument
  titlepage, % es wird eine Titelseite verwendet
  parskip=half, % Abstand zwischen Absätzen (halbe Zeile)
  headings=normal, % Größe der Überschriften verkleinern
  listof=totoc, % Verzeichnisse im Inhaltsverzeichnis aufführen
  bibliography=totoc, % Literaturverzeichnis im Inhaltsverzeichnis aufführen
  index=totoc, % Index im Inhaltsverzeichnis aufführen
  captions=tableheading, % Beschriftung von Tabellen unterhalb ausgeben
  final, % Status des Dokuments (final/draft)
  numbers=noenddot
]{scrreprt}

\renewcommand*{\chapterheadstartvskip}{\vspace{-1\baselineskip}} %Abstand vor Section
\setlength{\headsep}{30pt} %Abstand Kopfzeile Textkšrper


%Abstand vor und nach †berschrift einstellbar
%\makeatletter
%\def\@makeschapterhead#1{%
% \vspace*{5\p@}%       % 5 = Abstand vor chapter
%  {\parindent \z@ \raggedright
%    \normalfont
%    \interlinepenalty\@M
%    \huge \bfseries  #1\par\nobreak
%    \vskip 30\p@   % 30 = Abstand nach chapter
%  }}
%\makeatother


% Meta-Informationen -----------------------------------------------------------
%   Informationen über das Dokument, wie z.B. Titel, Autor, Matrikelnr. etc
%   werden in der Datei Meta.tex definiert und k\"onnen danach global
%   verwendet werden.
% ------------------------------------------------------------------------------
% Meta-Informationen -----------------------------------------------------------
%   Definition von globalen Parametern, die im gesamten Dokument verwendet
%   werden können (z.B auf dem Deckblatt etc.).
%
%   ACHTUNG: Wenn die Texte Umlaute oder ein Esszet enthalten, muss der folgende
%            Befehl bereits an dieser Stelle aktiviert werden:
%            \usepackage[latin1]{inputenc}
% ------------------------------------------------------------------------------
\newcommand{\titel}{Suchmaschinenmarketing am Beispiel eines Shops}
\newcommand{\untertitel}{\ }
\newcommand{\art}{Facharbeit}
\newcommand{\fachgebiet}{5. Semester}
\newcommand{\autor}{Christian Schramm}
\newcommand{\studienbereich}{Marketing}
\newcommand{\matrikelnr}{s3001102}
\newcommand{\erstgutachter}{Stefan Pabst}
\newcommand{\jahr}{2014}
\newcommand{\ort}{Dresden}
\newcommand{\logo}{ba_logo.jpg}


% ben\"otigte Packages -----------------------------------------------------------
%   LaTeX-Packages, die benötigt werden, sind in die Datei Packages.tex
%   "ausgelagert", um diese Vorlage möglichst übersichtlich zu halten.
% ------------------------------------------------------------------------------
% Anpassung des Seitenlayouts --------------------------------------------------
%   siehe Seitenstil.tex
% ------------------------------------------------------------------------------
\usepackage[
    automark, % Kapitelangaben in Kopfzeile automatisch erstellen
    headsepline, % Trennlinie unter Kopfzeile
    ilines % Trennlinie linksb�ndig ausrichten
]{scrpage2}

% Package f�r den Zeilenabstand
\usepackage[onehalfspacing]{setspace}
\recalctypearea

% Anpassung an Landessprache ---------------------------------------------------
\usepackage[ngerman]{babel}
\usepackage[babel,german=quotes]{csquotes}

% Umlaute ----------------------------------------------------------------------
%   Umlaute/Sonderzeichen wie ���� direkt im Quelltext verwenden (CodePage).
%   Erlaubt automatische Trennung von Worten mit Umlauten.
% ------------------------------------------------------------------------------
\usepackage[utf8]{inputenc}
\usepackage[T1]{fontenc}
\usepackage{textcomp} % Euro-Zeichen etc.
\usepackage{csquotes}

% 
\usepackage{scrhack}

% Schrift ----------------------------------------------------------------------
\usepackage{lmodern} % bessere Fonts
\usepackage{relsize} % Schriftgr��e relativ festlegen
%\usepackage{subscript}
\newcommand{\mysub}[1]{\raisebox{-0.5ex}{\scriptsize{#1}}}

% Grafiken ---------------------------------------------------------------------
% Einbinden von JPG-Grafiken erm�glichen
\usepackage[dvips,final]{graphicx}
% hier liegen die Bilder des Dokuments
\graphicspath{{Bilder/}}

% Befehle aus AMSTeX f�r mathematische Symbole z.B. \boldsymbol \mathbb --------
\usepackage{amsmath,amsfonts}

% f�r Index-Ausgabe mit \printindex --------------------------------------------
\usepackage{makeidx}

% Einfache Definition der Zeilenabst�nde und Seitenr�nder etc. -----------------
\usepackage{setspace}
\usepackage{geometry}

% Symbolverzeichnis ------------------------------------------------------------
%   Symbolverzeichnisse bequem erstellen. Beruht auf MakeIndex:
%     makeindex.exe %Name%.nlo -s nomencl.ist -o %Name%.nls
%   erzeugt dann das Verzeichnis. Dieser Befehl kann z.B. im TeXnicCenter
%   als Postprozessor eingetragen werden, damit er nicht st�ndig manuell
%   ausgef�hrt werden muss.
%   Die Definitionen sind ausgegliedert in die Datei "Glossar.tex".
% ------------------------------------------------------------------------------
\usepackage[printonlyused,footnote, nohyperlinks]{acronym}

% zum Umflie�en von Bildern ----------------------------------------------------
%\usepackage{floatflt}


% zum Einbinden von Programmcode -----------------------------------------------
\usepackage{listings}
\usepackage{xcolor} 
\definecolor{hellgelb}{rgb}{1,1,0.9}
\definecolor{colKeys}{rgb}{0,0,1}
\definecolor{colIdentifier}{rgb}{0,0,0}
\definecolor{colComments}{rgb}{1,0,0}
\definecolor{colString}{rgb}{0,0.5,0}
\lstset{
    float=hbp,
    basicstyle=\ttfamily\color{black}\small\smaller,
    keywordstyle=\color{colKeys},
    stringstyle=\color{colString},
    commentstyle=\color{colComments},
    columns=flexible,
    tabsize=2,
    frame=single,
    extendedchars=true,
    showspaces=false,
    showstringspaces=false,
    breaklines=true,
    backgroundcolor=\color{hellgelb},
    breakautoindent=true
}

% URL verlinken, lange URLs umbrechen etc. -------------------------------------
\usepackage{url}

% wichtig f�r korrekte Zitierweise ---------------------------------------------
\usepackage[square]{natbib}

%Farben definieren
\definecolor{darkblue}{rgb}{0,0,.5}

% PDF-Optionen -----------------------------------------------------------------

\usepackage[
    bookmarks,
    bookmarksopen=true,
    colorlinks=true,
% diese Farbdefinitionen zeichnen Links im PDF farblich aus
    linkcolor=blue, % einfache interne Verkn�pfungen
    anchorcolor=black,% Ankertext
    citecolor=blue, % Verweise auf Literaturverzeichniseintr�ge im Text
    filecolor=magenta, % Verkn�pfungen, die lokale Dateien �ffnen
    menucolor=red, % Acrobat-Men�punkte
    urlcolor=cyan, 
% diese Farbdefinitionen sollten f�r den Druck verwendet werden (alles schwarz)
    %linkcolor=black, % einfache interne Verkn�pfungen
    %anchorcolor=black, % Ankertext
    %citecolor=black, % Verweise auf Literaturverzeichniseintr�ge im Text
    %filecolor=black, % Verkn�pfungen, die lokale Dateien �ffnen
    %menucolor=black, % Acrobat-Men�punkte
    %urlcolor=black, 
    backref,
    plainpages=false, % zur korrekten Erstellung der Bookmarks
    pdfpagelabels, % zur korrekten Erstellung der Bookmarks
    hypertexnames=false, % zur korrekten Erstellung der Bookmarks
    linktocpage % Seitenzahlen anstatt Text im Inhaltsverzeichnis verlinken
]{hyperref}
% Befehle, die Umlaute ausgeben, f�hren zu Fehlern, wenn sie hyperref als Optionen �bergeben werden
\hypersetup{
    pdftitle={\titel \untertitel},
    pdfauthor={\autor},
    pdfcreator={\autor},
    pdfsubject={\titel \untertitel},
    pdfkeywords={\titel \untertitel},
}

% fortlaufendes Durchnummerieren der Fu�noten ----------------------------------
\usepackage{chngcntr}

% f�r lange Tabellen -----------------------------------------------------------
\usepackage{longtable}
\usepackage{colortbl}
\usepackage{array}
\usepackage{ragged2e}
\usepackage{lscape}

% Spaltendefinition rechtsb�ndig mit definierter Breite ------------------------
\newcolumntype{w}[1]{>{\raggedleft\hspace{0pt}}p{#1}}

% Formatierung von Listen �ndern -----------------------------------------------
\usepackage{paralist}

% bei der Definition eigener Befehle ben�tigt
\usepackage{ifthen}

% definiert u.a. die Befehle \todo und \listoftodos
\usepackage{todonotes}

% sorgt daf�r, dass Leerzeichen hinter parameterlosen Makros nicht als Makroendezeichen interpretiert werden
\usepackage{xspace}


% tables
\usepackage{booktabs}
\usepackage{multirow}


% Erstellung eines Index und Abkürzungsverzeichnisses aktivieren ---------------

\makeindex
%\makenomenclature

% Kopf- und Fußzeilen, Seitenränder etc. ---------------------------------------
% Zeilenabstand 1,5 Zeilen -----------------------------------------------------
\onehalfspacing

\sloppy

% Seitenr�nder -----------------------------------------------------------------
\setlength{\topskip}{\ht\strutbox} % behebt Warnung von geometry
\geometry{paper=a4paper,left=35mm,right=35mm,top=30mm}

% Kopf- und Fu�zeilen ----------------------------------------------------------
\pagestyle{scrheadings}
% Kopf- und Fu�zeile auch auf Kapitelanfangsseiten
\renewcommand*{\chapterpagestyle}{scrheadings} 
% Schriftform der Kopfzeile
\renewcommand{\headfont}{\normalfont}

% Kopfzeile
\ihead{\large{\textsc{\titel}}\\[1ex] \textit{\headmark}}
\chead{}
\ohead{\includegraphics[scale=0.3]{\logo}}
\setlength{\headheight}{21mm} % H�he der Kopfzeile

% Kopfzeile �ber den Text hinaus verbreitern
%\setheadwidth[0px]{textwithmarginpar} 
\setheadsepline[text]{0.4pt} % Trennlinie unter Kopfzeile

% Fu�zeile

\cfoot{}
\ofoot{\pagemark}

% sonstige typographische Einstellungen ----------------------------------------

% erzeugt ein wenig mehr Platz hinter einem Punkt
\frenchspacing 

% Schusterjungen und Hurenkinder vermeiden
\clubpenalty = 10000
\widowpenalty = 10000 
\displaywidowpenalty = 10000

% Quellcode-Ausgabe formatieren
\lstset{numbers=left, numberstyle=\tiny, numbersep=5pt, breaklines=true}
\lstset{emph={square}, emphstyle=\color{red}, emph={[2]root,base}, emphstyle={[2]\color{blue}}}

% Fu�noten fortlaufend durchnummerieren
\counterwithout{footnote}{chapter}

%Silbentrennung
\hyphenpenalty=10000
\exhyphenpenalty=10000


% eigene Definitionen für Silbentrennung
% Trennvorschl�ge im Text werden mit \" angegeben
% untrennbare W�rter und Ausnahmen von der normalen Trennung k�nnen in dieser
% Datei mittels \hyphenation definiert werden


% eigene LaTeX-Befehle
% Eigene Befehle und typographische Auszeichnungen f�r diese

% einfaches Wechseln der Schrift, z.B.: \changefont{cmss}{sbc}{n}
\newcommand{\changefont}[3]{\fontfamily{#1} \fontseries{#2} \fontshape{#3} \selectfont}

% Abk�rzungen mit korrektem Leerraum 
\newcommand{\ua}{\mbox{u.\,a.\ }}
\newcommand{\zB}{\mbox{z.\,B.\ }}
\newcommand{\dahe}{\mbox{d.\,h.\ }}
\newcommand{\Vgl}{Vgl.\ }
\newcommand{\bzw}{bzw.\ }
\newcommand{\evtl}{evtl.\ }

\newcommand{\abbildung}[1]{Abbildung~\ref{fig:#1}}

\newcommand{\bs}{$\backslash$}

% erzeugt ein Listenelement mit fetter �berschrift 
\newcommand{\itemd}[2]{\item{\textbf{#1}}\\{#2}}

% einige Befehle zum Zitieren --------------------------------------------------
\newcommand{\Zitat}[2][\empty]{\ifthenelse{\equal{#1}{\empty}}{\citep{#2}}{\citep[#1]{#2}}}

% zum Ausgeben von Autoren
\newcommand{\AutorName}[1]{\textsc{#1}}
\newcommand{\Autor}[1]{\AutorName{\citeauthor{#1}}}

% verschiedene Befehle um W�rter semantisch auszuzeichnen ----------------------
\newcommand{\NeuerBegriff}[1]{\textbf{#1}}
\newcommand{\Fachbegriff}[1]{\textit{#1}}

\newcommand{\Eingabe}[1]{\texttt{#1}}
\newcommand{\Code}[1]{\texttt{#1}}
\newcommand{\Datei}[1]{\texttt{#1}}

\newcommand{\Datentyp}[1]{\textsf{#1}}
\newcommand{\XMLElement}[1]{\textsf{#1}}
\newcommand{\Webservice}[1]{\textsf{#1}}


% Das eigentliche Dokument -----------------------------------------------------
%   Der eigentliche Inhalt des Dokuments beginnt hier. Die einzelnen Seiten
%   und Kapitel werden in eigene Dateien ausgelagert und hier nur inkludiert.
% ------------------------------------------------------------------------------
\begin{document}

% auch subsubsection nummerieren
\setcounter{secnumdepth}{3}
\setcounter{tocdepth}{3}

% Deckblatt und Abstract ohne Seitenzahl
\ofoot{}
\thispagestyle{plain}
\begin{titlepage}

\begin{center}

\huge{\textbf{\titel}}\\[1.5ex]
\LARGE{\textbf{\untertitel}}\\[6ex]
\LARGE{\textbf{\art}}\\[3ex]
\Large{\fachgebiet}\\[12ex]

\includegraphics[scale=0.5]{Logo_BA.jpg}\\[14ex]

\normalsize
\begin{tabular}{w{5.4cm}p{6cm}}\\
vorgelegt von:  & \quad \autor\\[1.2ex]
Studienbereich: & \quad \studienbereich\\[1.2ex]
Matrikelnummer: & \quad \matrikelnr\\[1.2ex]
Erstgutachter:  & \quad \erstgutachter\\[4ex]
\end{tabular}

\copyright\ \jahr\\[9ex]

\end{center}

%\singlespacing
%\small
%\noindent Dieses Werk einschlie�lich seiner Teile ist \textbf{urheberrechtlich gesch�tzt}. Jede Verwertung au�erhalb der engen Grenzen des Urheberrechtgesetzes ist ohne Zustimmung des Autors unzul�ssig und strafbar. Das gilt insbesondere f�r Vervielf�ltigungen, �bersetzungen, Mikroverfilmungen sowie die Einspeicherung und Verarbeitung in elektronischen Systemen.

\end{titlepage}




% !TEX root = Clean-Thesis.tex
%
\pdfbookmark[0]{Abstract}{Abstract}
\chapter*{Abstract}
\label{sec:abstract}
\vspace*{-10mm}

\blindtext

\vspace*{20mm}

{\usekomafont{chapter}Abstract (different language)}\label{sec:abstract-diff} \\

\blindtext



\ofoot{\pagemark}

% Seitennummerierung -----------------------------------------------------------
%   Vor dem Hauptteil werden die Seiten in großen römischen Ziffern
%   nummeriert.
% ------------------------------------------------------------------------------
\pagenumbering{gobble}
\tableofcontents % Inhaltsverzeichnis

\clearpage
\pagenumbering{arabic}
% die Inhaltskapitel werden in "Inhalt.tex" inkludiert -------------------------
% Hier k�nnen die einzelnen Kapitel inkludiert werden. Sie m�ssen in den 
% entsprechenden .TEX-Dateien vorliegen. Die Dateinamen k�nnen nat�rlich 
% angepasst werden.

\chapter{Einleitung}
\label{cha:Einleitung}

Dies wird die Einleitung für das Projekt.

Lorem ipsum dolor sit amet, consectetur adipisicing elit, sed do eiusmod tempor incididunt ut labore et dolore magna aliqua. Ut enim ad minim veniam, quis nostrud exercitation ullamco laboris nisi ut aliquip ex ea commodo consequat. Duis aute irure dolor in reprehenderit in voluptate velit esse cillum dolore eu fugiat nulla pariatur. Excepteur sint occaecat cupidatat non proident, sunt in culpa qui officia deserunt mollit anim id est laborum.
\chapter{Ziel und Anforderungen}
\label{cha:Allgemeines}

123

\section{Ziel des Projektes}
\label{sec:Ziel}

123

\chapter{Was ist Adware}
\label{cha:Adware}

\section{Section1}
\label{sec:Section1}
\chapter{Adsense}
\label{cha:Adsense}

123
\chapter{Umsetzung}
\label{cha:Umsetzung}

123





\chapter{Zusammenfassung}
\label{cha:Zusammenfassung}


 
 


123
%\include{Inhalt/AufzaehlungenTabellen}



\pagenumbering{Roman}

% Abkürzungsverzeichnis --------------------------------------------------------

%\chapter*
\addchap{Abk\"urzungsverzeichnis}
\begin{acronym}[laaaaaaaaaaang] %In der Eckigen Klammer muss das l�ngste Akronym stehen. Es wird dann die gesamte weitere Tabelle danach ausgerichtet
  
  \acro{ASP}{Active Server Pages}
  \acro{ERM}{Entity-Relationship-Modell}
  \acro{GUI}{Graphical User Interface}
  \acro{HTML}{Hypertext Markup Language}
  \acro{MSSQL}{Microsoft Structured Query Language}
  \acro{MVC}{Model View Controller}
  \acro{SQL}{Structured Query Language}
  
  
  


\end{acronym}

%\chapter*
\addchap{Quellen}

  http://de.wikipedia.org/wiki/Qualitätssicherung
  
  http://www.asp.net/mvc/mvc3
  
  http://www.software-architects.com/devblog/2008/03/15/Bug-Busters
  ---Test-Driven-Development-in-NET
  
  http://www.martinfowler.com/articles/continuousIntegration.html
  
  

\listoffigures % Abbildungsverzeichnis
\listoftables % Tabellenverzeichnis
%\renewcommand{\lstlistlistingname}{Verzeichnis der Listings}
%\lstlistoflistings % Listings-Verzeichnis

% arabische Seitenzahlen im Hauptteil ------------------------------------------





% Literaturverzeichnis ---------------------------------------------------------
%   Das Literaturverzeichnis wird aus der BibTeX-Datenbank "Bibliographie.bib"
%   erstellt.
% ------------------------------------------------------------------------------
%\bibliography{Bibliographie} % Aufruf: bibtex Diplomarbeit
% \bibliographystyle{natdin} % DIN-Stil des Literaturverzeichnisses

% Index ------------------------------------------------------------------------
%   Zum Erstellen eines Index, die folgende Zeile auskommentieren.
% ------------------------------------------------------------------------------
%\printindex



\addchap{Eidesstattliche Erklärung}
Ich, \autor, Matrikel-Nr.\ \matrikelnr, versichere hiermit, dass ich meinen Praxistransferbeleg  mit dem Thema
\begin{quote}
\textit{\titel} \textit{\untertitel}
\end{quote}
selbstständig verfasst und keine anderen als die angegebenen Quellen und Hilfsmittel benutzt habe, wobei ich alle wörtlichen und sinngemäßen Zitate als solche gekennzeichnet habe. Die Arbeit wurde bisher keiner anderen Prüfungsbehörde vorgelegt und auch nicht veröffentlicht.


\ort, den 
\newline

\rule[-0.2cm]{5cm}{0.5pt}

\textsc{\autor} 
 % Selbständigkeitserklärung

% Anhang -----------------------------------------------------------------------
%   Die Inhalte des Anhangs werden analog zu den Kapiteln inkludiert.
%   Dies geschieht in der Datei "Anhang.tex".
% ------------------------------------------------------------------------------
%\begin{appendix}
%    \clearpage
%    \pagenumbering{roman}
%    \chapter{Anhang}
%    \label{sec:Anhang}
%    % Rand der Aufzählungen in Tabellen anpassen
%    \setdefaultleftmargin{1em}{}{}{}{}{}
%    \section{DVD}
\label{sec:DVD}

\begin{description}
\item[Diplomarbeit]\begin{it}PDF-Datei und Anh"ange \end{it}
\begin{itemize}
\item{Diplomarbeit} 
\item{Dokument1}
\item{Dokument2} 
\item{Dokument3}
\item{Dokument4}
\item{Dokument5}
\item{Dokument6}
\item{Dokument7}
\end{itemize}
\item[Sonstiges]\begin{it}Zus"atzliche Informationonen zur Diplomarbeit \end{it}
\begin{itemize}
\item{Dokument1}
\item{Dokument2} 
\item{Dokument3}
\end{itemize}
\end{description}



\section{mkDoc2}
\label{sec:mkDoc2}
\vspace{4mm}
\includegraphics[width=0.3\textwidth]{ba_logo.jpg}

\section{mkDoc3}
\label{sec:mkDoc3}
\vspace{4mm}
\includegraphics[width=0.3\textwidth]{ba_logo.jpg}


%\end{appendix}


\end{document}
